\section{Introduction}
\IEEEPARstart{G}{raph} representation learning encodes graph-structured data into low-dimensional embeddings that capture topological and semantic information, underpinning relational modeling in domains such as social networks~\cite{sn_IJCAL,DBLP:conf/aaai/ManeK025,DBLP:conf/aaai/HuangLHZCZ25}, biochemistry~\cite{yang2023kerprint,xu2023seqcare}, and traffic systems~\cite{DBLP:conf/aaai/ZhangYXY0LZ25,fang2024efficient}. Within this context, Graph Neural Networks (GNNs)~\cite{GCN,Zhang2024InfiniteHorizonGF}, leveraging message-passing architectures, have become the dominant approach, enabling more expressive relational modeling than traditional node embedding methods~\cite{grover2016node2vec}. However, despite their effectiveness, GNNs often exhibit limited generalization to unseen graphs with substantially different topologies~\cite{zhao2024all}, a sensitivity that poses significant challenges for real-world applications where graph structures vary widely.



To overcome this limitation, retrieval-augmented methods for graph learning have recently emerged as a promising approach~\cite{lewis2021retrievalaugmented, JiangQXZZ00X0W24}, drawing inspiration from retrieval-augmented generation (RAG) techniques in NLP.
 RAG methods enhance GNNs by retrieving external subgraphs that are structurally or semantically relevant to the input, and integrating them into the learning pipeline to provide additional contextual information. Notably,  RAGRAPH~\cite{JiangQXZZ00X0W24} have demonstrated enhanced adaptability by incorporating retrieved graph fragments through message-passing prompt mechanisms. 
However, existing RAG methods operate primarily on nodes and edges, overlooking cycles that reflect latent higher-dimensional structures critical for capturing relational patterns and complex subgraph dependencies.%However, existing RAG methods focus mainly on low-dimensional elements (nodes and edges), overlooking higher-dimensional structures like cycles that are essential for retrieving meaningful substructures and reasoning over complex relations.

%However, existing RAG methods predominantly focus on low-dimensional elements—nodes and edges—failing to capture higher-dimensional topological characteristics, such as cycles, which are crucial for retrieving relevant substructures and enabling effective reasoning over complex relational dependencies.
\begin{figure}[t]
\centering\includegraphics[width=0.45\textwidth]{Pics/motivation_v2.pdf} 
    \caption{RAG-based graph learning with multi-dimensional topologies.}
    \label{fig:motivation}
  \end{figure}
  
In many real-world graphs (as illustrated in Fig.\ref{fig:motivation}), critical structural information arises not only from 0-dimensional nodes and 1-dimensional edges, but also from cyclic motifs, which form the boundaries of potential higher-dimensional cells. % within a cellular complex~\cite{DBLP:conf/nips/BodnarFOWLMB21}, a foundational notion in algebraic topology~\cite{hatcher_book}. 
These cycles capture recurring relational patterns and domain-specific structural dependencies, providing inductive biases crucial for accurate interpretation and retrieval. However, conventional GNN-RAG (Fig.\ref{fig:motivation}(a)) methods struggle to retrieve meaningful analogues when they fail to recognize or reason over such high-dimensional structures. In contrast, retrieving cyclic motifs from external graph corpora (Fig.\ref{fig:motivation}(b)) can provide structurally analogous examples that support inference in novel molecules, especially when query graphs exhibit global patterns that differ significantly from those during training.
 Moreover, high-dimensional reasoning captures complex dependencies beyond traditional pairwise message passing.

Despite the promise of topology-aware graph learning, leveraging multi-dimensional structures for retrieval-enhanced inference faces several key challenges. First, it is non-trivial to model graph substructures across multiple dimensions while preserving their hierarchical organization. Cell complexes~\cite{DBLP:conf/nips/BodnarFOWLMB21}, a fundamental construct in the theory of algebraic topology~\cite{hatcher_book},  provide a principled framework by organizing nodes, edges, and cycles into coherent units, where cycles serve as boundaries for higher-dimensional cells, capturing richer structural motifs than the original graph. Constructing such complexes consistently across heterogeneous graphs and organizing their multi-level components into a structured knowledge base is challenging. Second, effective reasoning requires propagating information across these multi-dimensional motifs. Relational dependencies often span nodes, edges, and higher-dimensional cells, demanding a mechanism that integrates messages across dimensions while preserving topological coherence—capabilities beyond conventional GNNs or standard retrieval approaches. Third, even with suitable structures retrieved, models must maintain their internal consistency; without explicit guidance, multi-dimensional information can collapse into pairwise interactions. A dedicated objective is therefore necessary to enforce coherent representations across cells and dimensions. 


%Moreover, reasoning over high-dimensional topological interactions enables the modeling of complex  dependencies—capturing rich structural semantics that traditional pairwise message passing schemes are fundamentally ill-equipped to express.
%In many real-world graphs, critical structural information emerges not only from 0-dimensional nodes and 1-dimensional edges, but also from higher-dimensional motifs—especially cycles—that together form regular cell complexes, a fundamental notion in algebraic topology~\cite{hatcher_book}. These higher-dimensional structures play a vital role across domains. For example, as shown in Fig.~\ref{fig:motivation}(b), in chemical compounds, for instance, cycles often correspond to functional molecular rings, whose presence and configuration directly determine chemical reactivity and biological activity.
%Such topological characteristics encode more than local interactions—they capture domain-specific inductive biases that are valuable for both interpretation and retrieval. GNN-RAG in Fig.~\ref{fig:motivation}(a) struggles to retrieve meaningful analogues when they cannot recognize or reason over high-dimensional topological features like cycles. 
%Retrieving cyclic motifs from external graph corpora can provide structural analogies that inform function prediction in novel molecules, particularly when the input graph exhibits global patterns not seen during training. Moreover, reasoning over high-dimensional topological interactions enables the modeling of multi-way dependencies that are difficult to express through traditional pairwise message passing.

In this work, we propose a novel \textbf{Re}trieved Cellular \textbf{T}opologies-\textbf{A}ugmented \textbf{G}raph Learning Framework (\textbf{ReTAG}) that leverages cellular complexes to model and retrieve multi-dimensional topology-aware subgraphs, termed cellular topologies. Specifically, ReTAG first lifts input graphs into cellular complexes to capture high-dimensional topological structures, enabling the extraction of multi-dimensional topology-aware substructures. These substructures are organized into a knowledge base that encodes rich topological and semantic information beyond traditional nodes and edges. During inference, ReTAG retrieves cellular topologies that align both topologically and semantically with the input graph, and integrates them through a multi-dimensional topological message-passing mechanism designed to propagate topological information effectively across different dimensions. Furthermore, a cellular topological contrastive learning module is introduced to enhance structural discrimination by enforcing feature consistency within each 2-cell, thereby capturing topological semantics beyond conventional pairwise interactions. %To further enhance learning, ReTAG incorporates topological supervision and optimization strategies that regularize the model and improve its generalization ability on unseen graphs. 
Extensive experiments demonstrate that ReTAG significantly outperforms existing methods across diverse graph scenarios.



Our contributions in this work are summarized as follows:
\begin{itemize}
    \item We highlight the critical role of higher-dimensional topological structures in graph learning.% and introduce the use of cellular complexes to model multi-dimensional topological interactions beyond traditional nodes and edges.

    \item We propose ReTAG that constructs a knowledge base of topology-aware substructures from lifted cellular complexes, and performs retrieval-augmented inference through multi-dimensional message passing to effectively integrate external topological knowledge.
    %ReTAG constructs a topology-aware knowledge base from lifted cellular complexes and performs retrieval-augmented inference through multi-dimensional message passing. Moreover, a cellular topological contrastive learning module is introduced to align features within each 2-cell, enabling the model to encode topological semantics beyond conventional edge-level relations.

    %Topology-Aware Learning: We design topological supervision and optimization techniques that regularize the learning process, enhancing model generalization on unseen graphs with diverse and complex structures.

    \item %Extensive experiments show ReTAG outperforms state-of-the-art baselines.
    Extensive experiments across multiple graph tasks demonstrate that ReTAG significantly outperforms state-of-the-art baselines.% by leveraging high-dimensional topological features, validating its robustness and effectiveness in modeling complex relational patterns.
\end{itemize}